\documentclass{chsmc}
\chsmcsetup{
	year = 2012,
	part = 1,
	type = B,
	show-answers = false,
}
\begin{document}

\section{填空题：本大题共 8 小题，每小题 8 分，共 64 分.}

\begin{enumerate}
	\item 对于集合 $\{x\mid a \leqslant x \leqslant b\}$，我们把 $b-a$ 称为它的长度.
		设集合 $A = \{x\mid a \leqslant x \leqslant a+1981\}$，
		$B = \{x\mid b-1014 \leqslant x \leqslant b\}$，且 $A$, $B$ 都是集合
		$U = \{x\mid 0 \leqslant x \leqslant 2012\}$ 的子集，则集合 $A \cap B$
		的长度的最小值是 \blankline
	\item 已知 $x > 0$, $y > 0$，且满足
		\[\begin{cases}
		\cos^2(\pi x) + 2\sin(\pi y) = 1, \\
		\sin(\pi x) + \sin(\pi y) = 0, \\
		x^2 - y^2 = 12,
		\end{cases}\]
		则有序数对 $(x,y) =$ \blankline
	\item 如图，设椭圆 $\dfrac{x^2}{a^2} + \dfrac{y^2}{b^2} = 1$（$a > b > 0$）的左、右焦点分别为
		$F_1$, $F_2$，过点 $F_2$ 的直线交椭圆于 $A(x_1,y_1)$, $B(x_2,y_2)$ 两点.
		若 $\triangle AF_1B$ 内切圆的面积为 $\pi$，且 $|y_1-y_2| = 4$，
		则椭圆的离心率为 \blankline
	\item 若关于 $x$ 的不等式组
		$$
		\begin{cases}
		x^3 + 3x^2 - x - 3 > 0, \\
		x^2 - 2ax - 1 \leqslant 0,
		\end{cases}
		$$
		（$a > 0$）的整数解有且只有一个，则 $a$ 的取值范围是 \blankline
	\item 在 $\triangle ABC$ 中，若 $\overrightarrow{AB}\cdot\overrightarrow{AC} = 7$，
		$|\overrightarrow{AB} - \overrightarrow{AC}| = 6$，则 $\triangle ABC$ 面积的最大值为 \blankline
	\item 如图，在长方体 $ABCD$-$A_1B_1C_1D_1$ 中，$AB = 4$，$BC = CC_1 = 2\sqrt{2}$，
		$M$ 是 $BC_1$ 的中点，$N$ 是 $MC_1$ 的中点. 若异面直线 $AN$ 与 $CM$ 所成的角为 $\theta$，
		距离为 $d$，则 $d\sin\theta =$ \blankline
	\item 设 $f(x)$ 是定义在 $\mathbf{R}$ 上的奇函数，且当 $x \geqslant 0$ 时，$f(x) = x^2$.
		若对任意的 $x \in [a,a+2]$，不等式 $f(x+a) \geqslant 2f(x)$ 恒成立，
		则实数 $a$ 的取值范围是 \blankline
	\item 一个均匀的正方体骰子的各面上分别标有数字 $1$, $2$, $\cdots$, $6$，
		每次投掷这样两个相同的骰子，规定向上的两个面上的数字之和为这次投掷的点数.
		那么，投掷 $3$ 次所得 $3$ 个点数之积能被 $14$ 整除的概率是 \blankline（用最简分数表示）
\end{enumerate}

\section{解答题：本大题共 3 个小题，满分 56 分. 解答应写出文字说明、证明过程或演算步骤.}

\begin{enumerate}[resume]
	\item (本题满分 16 分) 已知函数
		$f(x) = a\sin x - \dfrac{1}{2}\cos 2x + a - \dfrac{3}{a} + \dfrac{1}{2}$，
		$a \in \mathbf{R}$ 且 $a \neq 0$.
    \begin{enumerate}[(1)]
      \item 若对任意 $x \in \mathbf{R}$，都有 $f(x) \leqslant 0$，求 $a$ 的取值范围;
      \item 若 $a \geqslant 2$，且存在 $x \in \mathbf{R}$，使得 $f(x) \leqslant 0$，求 $a$ 的取值范围.
    \end{enumerate}
	\item (本题满分 20 分) 已知数列 $\{a_n\}$ 满足：当 $n = 2k - 1$（$k \in \mathbf{N}^*$）时，$a_n = n$；
		当 $n = 2k$（$k \in \mathbf{N}^*$）时，$a_n = a_k$. 记
		$T_n = a_1 + a_2 + a_3 + \cdots + a_{2^n-1} + a_{2^n}$，证明：对任意的 $n \in \mathbf{N}^*$，有
    \begin{enumerate}[(1)]
      \item $T_{n+1} = 4^n + T_n$;
      \item $\dfrac{1}{T_1} + \dfrac{1}{T_2} + \cdots + \dfrac{1}{T_n} < 1$.
    \end{enumerate}
	\item (本题满分 20 分) 已知抛物线 $y^2 = 2px$（$p > 0$），$A$, $B$ 是抛物线上不同于顶点 $O$ 的两个动点，
		记 $\angle AOB = \theta$（$\theta \neq 90^\circ$）. 若 $S_{\triangle AOB} = m\tan\theta$，
		试求当 $m$ 取得最小值时 $\tan\theta$ 的最大值.
\end{enumerate}

\end{document}
